\section{Formal coverage, open obligations, and ordinary Collatz}
\label{sec:formal-scope}

\subsection{What has actually been formalized}

The release contains three independent projects for the Lean~4 theorem
prover \cite{demoura-ullrich-2021}. It preserves the twelve
mathematical source modules of the initial research snapshot and incorporates
six additional signed-model modules from the subsequently merged audit. A
separate release bridge states the positive-seed capture result as an
if-and-only-if theorem. Thus there are nineteen mathematical modules:
eighteen source-preserved modules and one explicitly identified release
bridge. The original sources are unchanged. The source pins, byte comparisons,
and packaging transformations are recorded outside the minimal export;
\path{lean/README.md} explains the public layout.

The inverse project now includes the general signed physical model,
admissibility, height coordinates, cycle gain, and the three-family reduction,
in addition to its original three-step inverse results. It uses Lean 4.24.0,
as does the sprint project. The perturbed project uses Lean 4.19.0. None
requires mathlib. A clean verification builds each project in a new temporary
directory, checks the exact toolchain version, treats audit warnings as errors,
and compares the regenerated theorem inventory with the committed one.

\begin{center}\small
\begin{tabular}{lrrrr}\toprule
Project&Mathematical modules&Public&Private&Audited constants\\\midrule
\texttt{inverse}&11&110&37&556\\
\texttt{perturbed}&6&86&43&401\\
\texttt{sprint}&2&5&0&33\\\midrule
Total&19&201&80&990\\\bottomrule
\end{tabular}
\end{center}
The public and private columns count named source theorem declarations.
The last column also includes elaborator-generated theorem constants.
Two hundred public declarations come from preserved research sources; the
release bridge adds one. These counts are not independent-result counts
or percentages of the manuscript completed. The separate
\path{SevenAdversarial.lean} fixture checks twenty-one malformed-row or
scope examples and is not counted as a mathematical module.

The inventory in \path{lean/inventory.json} stores exact elaborated theorem
types, source locations, and transitive axiom dependencies. Across all 990
audited constants, the dependencies are drawn only from
\texttt{propext}, \texttt{Quot.sound}, and \texttt{Classical.choice}; no
\texttt{sorryAx} or custom axiom occurs. The complete audit includes unused
private and generated theorem helpers. An external independent Lean kernel
checker has not been run. The compile-and-audit result does not certify
unconnected Python, C++, plotting code, or the entire written manuscript.

The permitted axioms and transitive dependency inspection are documented
in the versioned Lean reference \cite[Sections~8.4--8.5]{lean-axioms-424}.
Axiom inspection does not decide whether the elaborated definitions and
theorem statement express the intended mathematics. The official
proof-validation guide treats that as a separate obligation and
distinguishes dependency auditing from proof replay and independent
kernel checking \cite{lean-validating-proofs}. The model comparisons
below address the statement-matching obligation; no additional checker
run is claimed here.

\paragraph{Directly matching results.}
The conservative cyclic-alternation assertion is proved in
\path{Collatz.Physical.conservative_admissible_iff_alternating}.
For arbitrary signed forcing, \texttt{admissible\_iff\_floor} gives the
binary-floor formulation of admissibility; the particular separated-negative-
impulse formulation in Theorem~\ref{thm:admissibility} is still presented
with its written proof. The source also proves binary-height existence
and uniqueness, the exact safe domain under the nonzero-word hypothesis,
physical/positive-coordinate conjugacy, its orbit and cycle consequences,
and admissibility under phase rotation. The primitive-period comparisons
in the manuscript are not separately formalized.

The all-depth inverse-lap and binary-zoom identities have direct Lean
support. So does the arithmetic return inequality $3^o<2^L$ used in
Proposition~\ref{prop:integer-cycle-attraction}. The latter theorem is
about integer arithmetic and safe returns; complex differentiation and
local analytic attraction are further written arguments.

The main new formal conclusion is
\path{Collatz.ReleaseBridge.global_capture_iff_three}:
universal capture of positive physical integers at all phases is
\emph{equivalent} to capture of the three complete B1, A3, N03 families.
It uses the committed signed theorem
\path{Collatz.SevenStep.global_capture_of_three} for the substantive
implication and proves the converse from positivity of the family launches.
The order of coordinates in Lean is phase then integer. Its physical word,
continued-zero convention, target cycle, and family parameter domains match
Section~\ref{sec:seven-three}. All 43 source Lean rows were independently
compared with the printed and Python rows, including the parameter shift
needed on zero-residue classes. The formal entry result is at most six
updates for positive seeds; the separate eleven-update safe-domain bound
and the scoped rank-optimality argument retain their stated evidence.

The family assumptions are explicit hypotheses, not trusted axioms. A
checked implication does not make those hypotheses true. The formal
statement therefore remains conditional on the same open mathematics as
the written theorem, while checking that particular implication in the signed physical model.

\paragraph{Related models still require care.}
The older perturbed project contains useful constituent lemmas with
different conventions. In particular, \texttt{PerturbedCollatz.step}
makes zero absorbing, and its orbit-lift theorem requires positivity
along the original prefix. Its period-seven example uses
$(-1,0,0,1,0,0,0)$, not the word in Section~\ref{sec:seven-step}.
Its escaping-basin theorem assumes an escaping target instead of
constructing one. The older progression-steering result uses a modulus
$n+1$ coprime to six and also requires $3^n\mid z-c$. It is not the
arbitrary-target all-progression theorem in this manuscript.

The sprint project supplies root inverse-height and packet identities;
those auxiliary declarations do not prove the residual-family hypotheses.
Likewise, the new signed parity-lift theorem supplies a positive odd affine
multiplier, but does not by itself formalize the entire parity-word/residue
bijection or identify it as $3$ raised to the number of odd inputs. Read the exact types before
transferring any statement. The claim-to-evidence map distinguishes a full
matching conclusion from a supporting identity and a related model.

\subsection{What remains to formalize and what requires new mathematics}

The signed model and the conditional three-family conclusion are now
available as formal foundations for further work. A next extension can
reuse them rather than introducing another convention for zero, phase,
or physical height. The remaining tasks have distinct sizes and proof
requirements.

For inverse growth, formalize the complete 24-step branch minima and
eligibility inequalities, then prove cross-depth noncollision, the
weighted stopping-frontier bound, and the finite selector's size and
multiplicity estimates. Correct arithmetic on a finite list does not
supply the list's exhaustive coverage or those universal counting
arguments automatically.

For the support bounds, a formal certificate interface must justify the
arbitrary-gap reduction, the complete signed tail closure, the finite
core/complement split, and every upper-bound enumeration. Replaying an
attaining witness or importing a recorded maximum is insufficient. The
current independent C++/Python replays remain computational evidence
outside Lean. For synchronization, the unfinished formal work includes
the protected drift, exponential tail, closed-reference conditioning,
union bound over all rotations, and primitive-period argument.

The continuous geometry would require real and complex analysis plus a
sound rational interval and disk interpreter. The proof obligations are
Taylor remainder bounds, outward rounding, complete interval coverage,
contraction, and the relation between disk capture and basin membership.
The numerical figure generator supplies none of these proof obligations.
Solovyev and Hales \cite[Sections~2--3]{solovyev-hales-2013} give a
relevant precedent: an external search produces Taylor/interval
certificates, which theorem-producing procedures then verify in
HOL~Light. For this release, a Lean interpreter would have to prove
the meaning and soundness of the actual rational certificate records.
The existing Python checker does not supply that formal interpreter.
The four-family theorem also remains a separate formalization task: its
control of returns beyond the short verification caps and its phase-correct
alternative-launch mergers are not formalized by the separate three-family
proof.

\begin{center}\small
\begin{tabular}{p{.25\textwidth}p{.33\textwidth}p{.31\textwidth}}\toprule
Target&Current evidence&Remaining obligation\\\midrule
Three-family global equivalence&Signed Lean proof and release iff bridge;
independent affine reconstruction&Prove the three whole-family capture
hypotheses to obtain unconditional capture\\[4pt]
Full foundation statements&Conservative classification, height/safety,
conjugacy and arithmetic cycle gain in Lean&Exact general separation
formulation, full parity-cylinder bijection, primitive-period statements\\[4pt]
Inverse growth in every progression&Written proof and exhaustive 24-step
certificate&Formal coverage, noncollision, weighted counting,
finite selectors and multiplicity\\[4pt]
Support bounds and synchronization&Complete written arguments and exact
support censuses&Formal upper-bound reduction/enumeration and the
probabilistic all-phase proof\\[4pt]
Continuous geometry&Rational disks, interval leaves and trace balls;
written analytic proofs&Sound analytic interval/disk interpreter and
formal contraction/basin arguments\\\bottomrule
\end{tabular}
\end{center}

Each milestone should end with a theorem whose complete type is compared
with the manuscript. Closing a formalization gap proves that a stated
argument has been checked by Lean. Closing B1, A3, and N03 requires new
mathematics about all parameters in those families; additional build logs
or declaration counts cannot perform that step.

\subsection{What would be needed for ordinary Collatz}

Let $T$ be the unforced shortcut map on positive integers. Its conjecture
asserts that every such orbit reaches 1. The results here concern either
a designed forcing or one of two particular fixed forced maps. The
coordinate change of Section~\ref{sec:foundations} gives a periodic
schedule of $T$ and $U$; it does not remove that schedule to give ordinary
$T$ iterations. No reduction from unforced capture to either set of
seven-step obligations has been established in this manuscript.

For example, a future transfer argument could specify an encoding $E$
of every positive unforced seed, prove a simulation relation between
unforced iterates and returns of a forced map, and show that forced
capture of every encoded state implies unforced capture. Such an argument
must prove the clock correspondence, guard coverage, safe domain,
existence of all required returns, and transfer of the terminal cycle.
Another possible route would prove a sufficiently strong exceptional-set
upper bound for exactly the forced map of Corollary~\ref{cor:exceptional-set},
then separately establish a sound transfer to the ordinary map. Neither
route is supplied by empirical similarity, inverse-basin abundance,
or a finite family of successful starting values.

The logarithmic-density result of Tao \cite{tao} is an example of a
powerful ordinary-Collatz theorem whose conclusion still differs from
universal capture: almost all orbits, in logarithmic density, attain a
bound given by any function tending to infinity. That statement is not
the exceptional-set power bound required here, and it concerns another
map. The distinction between an almost-everywhere conclusion and an
all-seeds conclusion remains substantive.

The three residual families in Theorem~\ref{thm:seven-three}, or
alternatively the four in Theorem~\ref{thm:seven-four}, are explicit
next targets within this research program. Progress would consist of an
exhaustive compatible reduction of one complete set of families, a
different complete invariant, or a
proof of their capture. Additional reductions on proper subprogressions
can be useful, but leave infinite complements unless an explicit coverage
argument closes them. The current work supplies several independently
checkable results and a concrete conditional reduction; it does not
establish the classical conjecture or claim to solve another named open
conjecture without a separate implication and literature check.
